% 特別研究II 論文(仮テンプレート)
% 研究室のテンプレートが手に入ったら、chapters/ の中身をそちらに流し込む。
% ビルド:thesis/ で latexmk main.tex → build/main.pdf
\documentclass[uplatex,dvipdfmx,a4paper,11pt,report,openany]{jsbook}

\usepackage{amsmath,amssymb,bm}
\usepackage{graphicx}
\usepackage{booktabs}
\usepackage{siunitx}
\usepackage{url}
\usepackage[dvipdfmx,hidelinks]{hyperref}
\usepackage{pxjahyper}

\graphicspath{{figures/}{../experiments/}}
\sisetup{detect-all}

% 表紙は先輩(2025年度 7307)の形式に合わせた
\newcommand{\thesistitle}{空間熱干渉を考慮した状態推定と\\動的熱マージンに基づくファン制御(仮題)}

\begin{document}
\begin{titlepage}
  \centering\Large
  \vspace*{40mm}
  令和 8 年度\quad 特別研究論文\par
  \vspace{30mm}
  \underline{題目\quad\parbox[t]{0.7\linewidth}{\centering\thesistitle}}\par
  \vspace{25mm}
  \underline{提出者\quad 7303\quad 鈴木\quad 黎映}\par
  \vspace{15mm}
  \underline{主査\quad 知念\quad 賢一\ 教授}\par
  \vfill
  \large 東京都立産業技術高等専門学校 専攻科\par
  \vspace{5mm}
  創造工学専攻\quad 情報工学コース\par
  \vspace*{20mm}
\end{titlepage}

\frontmatter
\chapter*{概要}
\addcontentsline{toc}{chapter}{概要}
% 最後に書く。先輩(20260115)の英文 Abstract の順番が参考になる:
% 背景 → 課題 → 従来研究の限界 → 提案 → 実装 → 数値つきの結果 → 効かなかった条件
（執筆中）

\tableofcontents
\listoffigures
\listoftables

\mainmatter
% !TEX root = ../main.tex
\chapter{はじめに}
% 約4ページ。最後に書く(目的の文言だけ早めに仮置き)

\section{研究背景}
% IoT・5G でエッジ処理が増え、ラズパイを束ねた小規模クラスタが安価な計算基盤として使われている

\section{エッジ環境における熱の問題}
% 専用空調のない居室に高密度で置くため排熱がこもりやすく、スロットリングで性能が落ちる

\section{空間熱干渉}
% 隣のノードの排熱で、自分は何もしていなくても温度が上がる。自ノードの負荷だけでは温度を説明できない
% ただし干渉の大きさは配置と囲いに強く依存する(4章:囲いなしではほぼ出ない)

\section{研究目的}
% 隣からの熱を推定しながら、必要なときだけファンを回し、スロットリングを避けつつファン稼働を減らす

\section{本論文の構成}

% !TEX root = ../main.tex
\chapter{関連技術}
% 約8〜10ページ。今ある材料で書ける。1節ずつ埋める

\section{シングルボードコンピュータとマイクロクラスタ}
% Raspberry Pi 4B の仕様(SoC BCM2711、4コア、最大クロック)。クラスタとしての利用例
% 設置形態:19インチラックのブレード型マウント、積み重ね型クラスタケースなど

\section{Raspberry Pi の熱特性}
% 温度センサ(/sys/class/thermal/thermal_zone0/temp、量子化は約0.5℃刻み)
% 約80℃からのソフト温度制限、DVFS でクロックが下がる仕組み
% vcgencmd get_throttled のビットの意味(例:0x80000 = ソフト温度制限が発生した)
% CPU ガバナ(ondemand / performance)

\section{冷却とファン}
% Raspberry Pi PoE+ HAT のファン(本研究の操作量)
% Linux からは冷却デバイス(type=pwm-fan、段階 0〜4)として見える
% OS 標準の制御(thermal_zone0 の step_wise、40/45/50/55℃ で段階的に回る)と、
% それを user_space に切り替えて自前で制御する方法

\section{熱RCモデル}
% 熱容量 C と熱抵抗 R:C dT/dt = -(T - T_amb)/R + P、時定数 tau = RC
% 離散化(5章の導出の準備)

\section{状態空間モデルとカルマンフィルタ}
% 予測と更新、カルマンゲイン、誤差共分散 P が推定の不確かさを表すこと
% 多段先予測での P の伝播

\section{計測・負荷生成の技術}
% vcgencmd、/proc/stat による CPU 使用率(区間平均)、MQTT、stress-ng、時刻同期(NTP)
% SwitchBot 温湿度計(室温の記録)

% !TEX root = ../main.tex
\chapter{関連研究}
% 約5ページ。IPSJ 予稿の参考文献を土台に広げる

\section{データセンタの熱再循環と熱を考慮した配置}
% Tang, Moore ら

\section{予測に基づく熱管理}
% Zanini, Skadron ら:熱RCモデル・MPC。モデルが正確な前提で、想定外の外乱に弱い

\section{不確かさを考慮したマージン設計}
% Mesbah、坂本ら:過去の誤差の統計でマージンを決める

\section{エッジ・SBC環境の熱管理}
% EdgeEngine など

\section{本研究の位置づけ}
% 比較表:干渉を扱うか/不確かさを使うか/実機で制御するか/ファン稼働を評価するか

% !TEX root = ../main.tex
\chapter{予備調査}
% 約5〜6ページ。データはそろっている

\section{先行発表(IPSJ)の実験と結果}
% 2台、タオルで囲った状態(CPU 面を向かい合わせ)、ファンなし
% 片側負荷で隣が約+12℃(data/legacy/expA, expB)、両側負荷では約84℃で頭打ち(expC)

\section{先行発表の課題}
% legacy/ipsj2026/linux_pc_sotuken/README.md と experiments/20261112_legacy_kf_check/README.md
% ・Δ秒先予測に未来の実測値(CPU使用率・隣接温度)を使っていた
% ・同定と評価が同じデータ。同じデータで設定を変えながら評価した
% ・CPU 使用率が瞬時値で記録され、負荷の係数 b が負(物理的に不自然)
% ・制御は実機に反映されていない(オフライン評価)。miss_count の定義の問題
% ・実験は約20分×1回、室温未記録、囲い(タオル)の再現性がない

\section{囲いなし・ラックでの熱干渉}
% experiments/20261113_openair_interference:囲いなしでは隣も離れたノードも +0.5℃ で差なし
% experiments/20261113_rack_interference:(結果待ち)

\section{本研究の方針}
% 干渉が出る条件と出ない条件を両方示す/干渉を受ける側も評価する/実際にファンを動かす/
% 室温を記録し、条件と手順を固定して複数回行う

\chapter{提案手法}
% 約8〜10ページ。式は「なぜこうなるか」を必ず書く。下書きは docs/parameters.md

\section{全体の流れ}
% 観測 → 推定 → Δ秒先予測 → マージン決定 → ファン ON/OFF(ブロック図)

\section{自ノードの熱モデルの導出}
% RC モデル → P = P0 + k u → 厳密な離散化:a = exp(-dt/tau), b = (1-a) R k, d = (1-a)(T_amb + R P0)

\section{空間熱干渉のモデル化}
% H_{t+1} = rho H_t + (1-rho) z_t、z_t = 隣の温度 - 基準温度。rho = exp(-dt/tau_H)

\section{ファンの効果のモデル化}
% 案A:ON/OFF でパラメータを切り替える/案B:+ e f_t。同定データで比べて決める

\section{パラメータの同定}
% 同定専用のステップ実験(experiments/*_step_identify)。物理的な妥当性(0<a<1, b>0)の確認

\section{カルマンフィルタによる推定}
% T と H の同時推定、適応プロセスノイズ qH = qH0 + beta |Δz|

\section{Δ秒先の予測と不確かさ}
% 未来の入力は「今の値が続く」と仮定(IPSJ 版との違い)。σ = sqrt(P_TT)
% Δ の決め方:ファンの応答時間 < Δ < τ

\section{動的マージン}
% M = alpha sigma + delta、Tsafe = Tlimit - M。alpha = 2 の根拠(片側 2.3%)

\section{ファン制御則}
% 予測温度 > Tsafe で ON、最小 ON 時間とヒステリシス

% !TEX root = ../main.tex
\chapter{実装}
% 約5ページ。もう実装済みなので書ける

\section{システム構成}\label{sec:impl-system}
図\ref{fig:system}に本研究のシステム構成を示す。本システムは、3台のノードと、計測データの収集と制御を行う制御用 PC 1台からなる。

\begin{figure}[tb]
  \centering
  \fbox{\parbox{0.8\linewidth}{\centering\vspace{20mm}システム構成図(作成中)\vspace{20mm}}}
  \caption{システム構成}
  \label{fig:system}
\end{figure}

ノードには Raspi を3台用いた。3台はラックの上段に約 XX\,mm 間隔で横一列に設置し、左から raspi1、raspi2、raspi3 とした。各機器の仕様を表\ref{tab:env}に示す。各ノードへの給電は PoE スイッチから行った。各ノードには Raspberry Pi PoE+ HAT を装着し、そのファンを冷却に用いる。このファンはソフトウェアから ON/OFF でき、本研究の制御における操作量である。


\begin{table}[tb]
  \centering
  \caption{実装環境}
  \label{tab:env}
  \begin{tabular}{lll}
    \toprule
    & ノード(3台) & 制御用 PC(1台) \\
    \midrule
    機種 & Raspberry Pi 4 Model B & Lenovo IdeaPad 5 14ALC05 \\
    CPU & Cortex-A72 4 コア & AMD Ryzen 5 5500U \\
    メモリ & 4\,GB & 8\,GB \\
    OS & Ubuntu 24.04.3 LTS & Ubuntu 24.04.4 LTS \\
    冷却 & PoE+ HAT のファン & --- \\
    給電 & PoE & AC アダプタ \\
    Python & 3.12.3 & 3.12.3 \\
    MQTT & paho-mqtt 1.6.1 & mosquitto 2.0.18,paho-mqtt 2.1.0 \\
    負荷生成 & stress-ng 0.17.06 & --- \\
    数値計算 & --- & NumPy 2.4.1,pandas 2.3.3 \\
    \bottomrule
  \end{tabular}
\end{table}


制御用 PC では MQTT ブローカーとして mosquitto を動作させた。
各ノードは1秒ごとに計測値を MQTT で制御用 PC へ送信し、制御用 PC はファンの指令と、実験時には負荷の指令を MQTT で各ノードへ送信する。データの収集と記録、温度の推定、ファンの ON/OFF の判断は、すべて制御用 PC で行う。

推定と制御をノード上ではなく別の PC で行った理由は2つある。1つ目は、ノード間の条件をそろえるためである。特定のノードで制御プログラムを動かすと、そのノードだけ処理による発熱が加わり、測定したい温度に影響する。2つ目は、制御の遅れを避けるためである。高負荷で温度が上がったノードで制御プログラムを動かすと、ファンを回す判断が最も必要なときに処理が遅れるおそれがある。なお、制御用 PC が停止した場合の対策は\ref{sec:impl-fan}節で述べる。




\section{データ収集}
% scripts/node/node_agent.py:1秒ごとに温度、CPU使用率(/proc/stat の差分)、クロック、get_throttled、
% ファン状態、PoE 電流。systemd サービスとして動かす(scripts/node/install_node.sh)
% scripts/ambient_logger.py:SwitchBot 温湿度計(ラックの上・左・右)とハブ(室温)を1分ごと

\section{推定・制御プログラム}
% src/etc/(model.py, kalman.py)

\section{ファンの駆動}\label{sec:impl-fan}
% PoE+ HAT のファンを冷却デバイス経由で ON/OFF。実験中は OS の制御を止め、終了時に戻す
% 安全機能:82℃ 以上、または指令が10秒途絶えたら強制 ON

\section{負荷生成と実験の自動化}
% scripts/run_experiment.py:シナリオ(JSON)、開始前に温度が落ち着くまで待つ、切り替え時刻を events.csv に記録

% !TEX root = ../main.tex
\chapter{評価実験}
% 約10ページ。予備審査では予備結果まで

\section{実験条件}
\section{熱干渉の特性}
\section{予測精度}
% Self-only / 坂本型 / H なし KF / 提案。負荷側 A と隣 B を分けて
\section{予測ホライズンの感度分析}
\section{ファン制御}
% ファンなし / 常時ON / しきい値 / Self-only+固定 / KF+固定 / 提案。ファン稼働率と最高温度のトレードオフ

\chapter{考察}
\section{熱干渉が効く条件}
% 配置(熱が昇る方向に隣がいるか)と囲いで決まる
\section{外部熱影響の推定の効果}
\section{動的マージンとファン稼働の関係}
\section{パラメータの物理的な妥当性}
\section{本研究の限界}
% 室温の変動、台数、ON/OFF のみ、恒温室なし

\chapter{おわりに}
\section{まとめ}
\section{今後の課題}


\backmatter
\chapter*{謝辞}
\addcontentsline{toc}{chapter}{謝辞}
% TODO

\nocite{*} % 下書きの間は全文献を載せる。本文で \cite するようになったら消す
\bibliographystyle{junsrt}
\bibliography{references}
\end{document}
